\chapter{Architettura dati dei profili}
\label{ch:architettura_dati_profili}

I profili demo (\texttt{small}, \texttt{medium}, \texttt{big},
\texttt{huge}, \texttt{superhuge}, \texttt{mega}) sono il banco di
prova quotidiano del solver: l'idea \`e di sottoporre l'engine a
sei taglie progressive di scuola (10\,$\to$\,100 classi) gi\`a corredate
di un realistico ventaglio di vincoli, cos\`i che chi importa il
profilo dalla dashboard si ritrovi davanti uno scenario \emph{stress}
che esercita ogni famiglia di vincolo gestita dal solver. Per anni
i profili erano due pickle slim --- \texttt{school\_<profilo>.pkl}
con l'anagrafica e \texttt{profs\_<profilo>.pkl} con le cattedre ---
e l'import generava al volo le aule, gli studenti e le ore di
lavoro tramite gli helper di mock. \`E un modello che ha funzionato
finch\'e i vincoli erano pochi, ma quando le tabelle di constraint
hanno superato la dozzina si \`e fatto pesante: ogni import doveva
ricostruire da zero tutto ci\`o che il pickle non sapeva esprimere,
e in pi\`u cancellava eventuali fixture lasciate dall'utente nelle
tabelle vincoli. Da questo capitolo in poi la fonte di verit\`a per
i profili \`e un singolo SQLite per profilo, costruito una volta
sola da \texttt{engine/scripts/build\_profile\_db.py} con seed
deterministico e poi importato \emph{verbatim} dal backend.

\section{Anatomia di un profilo SQLite}

Ogni file \texttt{engine/scripts/data/<profilo>/<profilo>.sqlite}
contiene esattamente lo stesso schema di
\texttt{webui/data/timetable.db}, costruito tramite
\texttt{Base.metadata.create\_all} sullo stesso \texttt{models.py}
del runtime. Importare lo SQLite vuol dire copiare riga per riga
le tabelle di interesse nella DB live, con la sola accortezza di
proiettare la copia sull'intersezione delle colonne presenti in
entrambi i database --- una salvaguardia che rende il pipeline
resiliente quando la DB live ha gi\`a ricevuto una migration che lo
SQLite del profilo non ha ancora visto, o viceversa. Per orientarsi
nel contenuto del file basta aprirlo con \texttt{sqlite3} e
contare le righe nelle famiglie di tabelle: anagrafica, vincoli,
plessi, calendario, soluzione.

\begin{figure}[h]
\centering
\begin{tikzpicture}[node distance=0.55cm and 0.7cm,
  every node/.style={draw, rounded corners=4pt, very thick,
                     minimum width=4.6cm, minimum height=0.95cm,
                     align=center, font=\small}]
  \node[fill=cBox!30] (an) {Anagrafica\\
    \textit{subjects, teachers,}\\
    \textit{school\_classes, assignments}};
  \node[fill=cBox!30, right=of an] (rm) {Aule + plessi\\
    \textit{classrooms, plessi,}\\
    \textit{plesso\_*\_rules/policies}};
  \node[fill=cBox!30, below=of an] (cs) {14 tabelle vincoli\\
    \textit{teacher/class/room\_unavail,}\\
    \textit{coteach\_groups, gen.\,DSL}};
  \node[fill=cBox!30, below=of rm] (oh) {Calendario settimanale\\
    \textit{working\_days,}\\
    \textit{working\_hour\_slots}};
  \node[fill=cBox!30, below=2cm of cs.south] (sl) {Soluzione (opzionale)\\
    \textit{solutions, lessons}};
  \draw[-{Stealth[length=2.5mm]}, thick] (an) -- (cs);
  \draw[-{Stealth[length=2.5mm]}, thick] (rm) -- (oh);
  \draw[-{Stealth[length=2.5mm]}, thick] (an) -- (rm);
  \draw[-{Stealth[length=2.5mm]}, thick] (cs) -- (sl);
  \draw[-{Stealth[length=2.5mm]}, thick] (oh) -- (sl);
\end{tikzpicture}
\caption{Anatomia di \texttt{<profilo>.sqlite}: cinque blocchi
   tematici, tutti popolati da \texttt{build\_profile\_db.py} con
   seed deterministico.}
\label{fig:profilo_sqlite}
\end{figure}

Il file \`e autosufficiente nel senso che ogni \emph{constraint} che
verrebbe altrimenti sintetizzato a runtime --- una \texttt{teacher\_unavailability}
con il giorno bloccato HARD, una \texttt{coteach\_groups} che
forza tre ore di compresenza in laboratorio, una \texttt{plesso\_commuting\_rules}
che impedisce a un docente di fare i primi 60 minuti in succursale
e i secondi a 800 metri di distanza --- vive direttamente come
riga della tabella corrispondente, con stato HARD o SOFT esplicito
e la stessa shape che il modello SQLAlchemy si aspetta in produzione.
Lo script di build cura anche un \texttt{manifest.json} fianco a fianco
col file SQLite: contiene il seed, l'etichetta dello schema orario
del profilo, e i conteggi per ogni tabella, in modo che la sezione
``Stress per profilo'' del capitolo Benchmark si rigeneri da zero
con un \texttt{python -m engine.scripts.build\_profile\_db --all}.

\section{Le 14 tabelle di vincolo}

L'aggettivo ``quattordici'' \`e meno ostico di quanto suggerirebbe il
solo elenco, perch\'e le tabelle si distribuiscono naturalmente su
tre assi --- \emph{chi}, \emph{quando}, \emph{come} --- e il profilo
le tocca tutte con almeno una riga. Sull'asse \emph{chi} la prima
trinit\`a riguarda i docenti: \texttt{teacher\_unavailability}
(stato HARD/SOFT per cella giorno-ora),
\texttt{teacher\_mandatory\_free\_days} (HARD: 0 ore in quel
giorno, tipico per part-time) e
\texttt{teacher\_compatible\_classes} (allow-list HARD per cattedre
ristrette). A queste si aggiungono le due tabelle di preferenza
con tassonomia a cinque stati,
\texttt{teacher\_class\_preferences} e
\texttt{teacher\_curriculum\_preferences}, che il solver di Phase A
legge come HARD nei rami estremi (\texttt{enforced} / \texttt{forbidden})
e SOFT nei rami intermedi.

L'asse \emph{quando} \`e dominato dalle tre tabelle di disponibilit\`a
non centrate sul docente: \texttt{class\_unavailability},
\texttt{classroom\_unavailability} e
\texttt{logical\_unavailabilities}, quest'ultima con DNF
parametrica per esprimere espressioni come ``il prof X \`e libero
solo nel pomeriggio del marted\`i e gioved\`i''. La quarta tabella
dello stesso asse, \texttt{curriculum\_logical\_constraints}, vive
fuori dal rispetto temporale puro per occuparsi di vincoli
trasversali al curriculum --- ``Matematica e Italiano non lo stesso
giorno della prima'' \`e l'esempio di stress che il profilo small
scrive a costruzione.

L'asse \emph{come} raggruppa le strutture di compresenza,
\texttt{coteach\_groups}, e i due moduli plesso-aware:
\texttt{plesso\_commuting\_rules} (regole tra coppie di plessi,
HARD o SOFT con peso) e \texttt{plesso\_entity\_policies}
(politiche per docente o classe del tipo
\texttt{single\_plesso\_per\_day} oppure \texttt{single\_plesso\_total}).
Manca all'appello la quattordicesima --- \texttt{general\_constraints} ---
che tecnicamente non \`e una tabella di vincolo ``classico'' ma
contiene espressioni nel DSL generico (vedi Capitolo
\ref{ch:dsl_generico_per_i_vincoli}) che il post-processor valuta
sulla soluzione attiva. Il profilo small ne scrive cinque,
inclusa \texttt{forall l in lessons where l.subject==Educazione Fisica:
l.classroom.kind==palestra} che esercita la regola HARD
``EduFisica solo in palestra''.

\section{Il vincolo di capacit\`a}

La novit\`a di pi\`u largo impatto introdotta insieme ai profili
SQLite \`e il vincolo HARD di capacit\`a: per ogni \texttt{Lesson}
deve valere \texttt{lesson.classroom.capacity >=
lesson.school\_class.n\_students}. La regola \`e implementata
nativamente in \texttt{engine/classroom\_assignment.py}, dove la
funzione \texttt{\_can\_host} che decide l'eleggibilit\`a aula-lezione
filtra ogni candidato troppo piccolo prima ancora di costruire il
modello CP-SAT. La stessa regola \`e disponibile come pragma DSL
\texttt{classroom\_capacity\_ok()} nel compilatore
\texttt{engine/dsl\_to\_cpsat.py} per i solver che costruiscono
uno spazio di decisione congiunto slot+room (la
MonolithicSolver con un \texttt{classroom\_for\_slot} noto): in
quel caso il pragma scorre lo spazio dei slot, recupera la
capacit\`a dell'aula assegnata a ciascun \texttt{(class, day, hour)}
e forza a zero la \texttt{BoolVar} corrispondente quando la
capienza \`e insufficiente. Il pragma diventa un no-op con
diagnostica ``classroom data empty'' nei contesti che non hanno
ancora le aule assegnate, perch\'e in quel caso \`e il solver di
Phase~C a rendere conto del vincolo.

Il dato di ingresso \`e popolato per ogni profilo dallo script di
build: tutte le classi prendono un \texttt{n\_students} con
distribuzione gaussiana troncata su $[18, 30]$ (centrata su 24,
$\sigma=3.5$), e una o due classi vengono volutamente forzate a
combaciare esattamente con la capienza dell'aula standard pi\`u
grande del profilo. Il risultato \`e un vincolo \emph{stringente}:
la classe pu\`o essere ospitata solo in quella precisa aula, con
tutte le altre messe sotto pressione ad accomodare gli orari
intorno. L'utente vede a colpo d'occhio l'impatto del vincolo nel
tab \texttt{/classes} (colonna ``N. studenti'' editabile inline,
range 5--50) e in \texttt{/classrooms} (colonna ``Capienza''
editabile inline, range 5--100); il backend \`e autoritativo,
ma le finestre di creazione lezione (\texttt{AddLessonModal} e
\texttt{AddEventModal}) mostrano un avviso rosso a tempo di
selezione quando la coppia (classe, aula) viola la regola.

\section{Aule speciali e pragma di tipologia}

Accanto al vincolo di capacit\`a vive il vincolo di tipologia:
\texttt{Subject.required\_kind} \`e il modo a singola riga di
esprimere ``Educazione Fisica solo in palestra'' o ``Fisica solo
in laboratorio fisica'', senza dover enumerare tutte le altre aule
con preferenze \texttt{forbidden}. Il valore \`e una stringa
nello stesso vocabolario di \texttt{Classroom.kind} (\texttt{standard},
\texttt{palestra}, \texttt{lab\_fisica}, \texttt{lab\_chimica},
\texttt{lab\_informatica}, \texttt{lab\_linguistico},
\texttt{biblioteca}, \texttt{aula\_speciale}); \texttt{NULL}
significa ``qualsiasi''. La funzione \texttt{\_can\_host} legge
il campo dalla \texttt{Lesson} (proiettato nel dict che il
backend produce in \texttt{lessons\_for\_classroom\_step}) e
rifiuta ogni aula con \texttt{kind} diverso. Lo stesso vincolo
viene in pi\`u espresso come pragma DSL nelle
\texttt{general\_constraints} dei profili, perch\'e l'utente che
apre il tab \texttt{/general-constraints} possa modificare la
regola direttamente da UI senza toccare il modello.

\section{Schemi orari differenziati per profilo}

Lo stress di un profilo non \`e solo numero di righe: \`e anche
varianza nello \emph{schema} con cui le righe si dispongono nello
spazio settimanale. La generazione del calendario, perci\`o, non \`e
un copia-incolla: ogni profilo riceve la sua versione di
\texttt{working\_days} + \texttt{working\_hour\_slots} pensata per
mettere in difficolt\`a un aspetto diverso del solver.

\begin{table}[h]
\centering
\caption{Schemi orari per profilo (cfr. \texttt{engine/scripts/build\_profile\_db.py}, funzione \texttt{working\_schedule}).}
\label{tab:profili_orari}
\small
\begin{tabular}{lp{8.5cm}}
\toprule
\textbf{Profilo} & \textbf{Schema settimanale} \\
\midrule
\texttt{small}     & 8:00--14:00 lun--sab, 6 slot orari da 60 min. \\
\texttt{medium}    & 8:00--14:00 lun--ven, 8:00--12:00 sab. \\
\texttt{big}       & 8:00--14:00 lun--ven, 8:00--13:00 sab. \\
\texttt{huge}      & 7:55--13:55 lun--sab, slot da 55 min con
                     5 min di pausa baked-in. \\
\texttt{superhuge} & 8:00--13:00 lun--ven + 14:00--17:00 mar/gio
                     (mattina + pomeriggio bisettimanale). \\
\texttt{mega}      & 8:00--14:00 lun--ven + 8:00--13:00 sab +
                     14:00--17:00 mer (un pomeriggio fisso). \\
\bottomrule
\end{tabular}
\end{table}

Il dato chiave \`e che ogni schema mantiene la corrispondenza con
i numeri di giorno legacy 1--6 e di ora 8--13 attraverso i campi
\texttt{legacy\_day\_number} e \texttt{legacy\_hour\_number} delle
righe \texttt{working\_days} / \texttt{working\_hour\_slots}, in
modo che le tabelle di \emph{unavailability} che fanno riferimento
a quei numeri sopravvivano a un eventuale rebuild dell'orario senza
dover migrare i dati uno per uno. Il profilo \texttt{huge} \`e
un caso interessante: i suoi slot da 55 minuti restano con
\texttt{legacy\_hour\_number} 8--13, e il blocco di 5 minuti tra
una lezione e la successiva non \`e modellato come slot intermedio
ma come margine implicito tra \texttt{end\_time} di slot $i$ e
\texttt{start\_time} di slot $i+1$. Cos\`i il solver continua a
ragionare in slot, ma l'utente che apre la calendar view vede
gli orari con la pausa.

\section[Provalo subito]{Provalo subito\protect\footnote{La sezione
prende il nome dalla rubrica omonima del README -- riproducibilit\`a
``a portata di copia-incolla''.}}

Per rigenerare lo SQLite di un singolo profilo bastano due
righe di shell --- la prima invoca lo script con il nome del
profilo, la seconda apre il file con \texttt{sqlite3} per un
controllo a campione delle tabelle:

\begin{lstlisting}[style=shell]
python -m engine.scripts.build_profile_db small
sqlite3 engine/scripts/data/small/small.sqlite \
    "select count(*) from general_constraints;"
\end{lstlisting}

\noindent Il primo comando produce
\texttt{engine/scripts/data/small/small.sqlite} insieme al
\texttt{manifest.json} e a \texttt{PICKLE\_DEPRECATED.md}; il
secondo restituisce il numero di righe nella tabella delle
\texttt{general\_constraints} --- cinque, in tutti i profili. Per
ricostruire tutti e sei i profili in serie (in subprocess separati
perch\'e ogni profilo va in un suo file SQLite e SQLAlchemy non
accetta volentieri di rebindare l'engine a runtime) basta sostituire
il primo comando con \texttt{python -m engine.scripts.build\_profile\_db --all};
lo script aggrega in coda i singoli \texttt{manifest.json} dentro
\texttt{engine/scripts/data/profiles\_manifest.json}, che la sezione
``Stress per profilo'' del capitolo Benchmark legge per popolare
la tabella riepilogativa.

\section{Compatibilit\`a con la pipeline pickle}

I file \texttt{school\_<profilo>.pkl} e
\texttt{profs\_<profilo>.pkl} restano nel repository per finalit\`a
storiche --- audit, debug di import legacy, comparativa con i run
del trimestre scorso --- ma non sono pi\`u letti come default
dall'\texttt{import\_engine\_profile}: la routine cerca prima
\texttt{<profilo>.sqlite} nella stessa cartella e usa il pickle
solo se lo SQLite manca. Una nota \texttt{PICKLE\_DEPRECATED.md}
accanto a ogni pickle elenca le tabelle e i conteggi attuali per
quel profilo, cos\`i che chi apre la cartella non si chieda perch\'e
l'import si comporta diversamente da quanto memorizza il
\texttt{git log}: la ragione \`e che lo SQLite \`e ora la fonte
autoritativa, ed \`e l'unico file che il backend si aspetta di
trovare. La porta verso il vecchio modo resta aperta, ma il
percorso default cambia.
